\chapter{Modelamiento Conceptual Entidad-Relaci\'on (DER/EER) y Cardinalidad}
\label{cap:modelamiento_conceptual_der}

\section{Marco Conceptual y Te\'orico (Curr\'iculo Universitario)}

\subsection{El Modelo Entidad-Relaci\'on Extendido (EER)}
El modelamiento conceptual traduce los requerimientos informales de negocio en especificaciones abstractas no dependientes de ning\'un motor f\'isico espec\'ifico.

\begin{itemize}
    \item \textbf{Entidades Fuertes y D\'ebiles:} Las entidades fuertes poseen existencia independiente y clave propia; las entidades d\'ebiles dependen existencialmente de otra entidad y se identifican mediante una clave parcial m\'as la clave for\'anea del propietario.
    \item \textbf{Tipos de Atributos:}
    \begin{itemize}
        \item \textit{Simples / At\'omicos:} No divisibles (ej. \texttt{precio}).
        \item \textit{Compuestos:} Jer\'arquicos (ej. \texttt{Direccion: Calle, Ciudad, CP}).
        \item \textit{Multivaluados:} Colecciones de valores para una misma instancia (ej. \texttt{telefonos[]}).
        \item \textit{Derivados:} Calculados en tiempo de ejecuci\'on (ej. \texttt{edad = CURRENT\_DATE - fecha\_nacimiento}).
    \end{itemize}
    \item \textbf{Cardinalidad y Participaci\'on:} Ratios estructurales $(1:1, 1:N, N:M)$ y participaci\'on total (obligatoria) o parcial (opcional).
\end{itemize}

---

\section{Casos de Estudio: Mapeo de Relaciones N:M a Tablas Puente}

Cuando existe una relaci\'on Muchos a Muchos entre entidades de negocio (por ejemplo, \textbf{\'ordenes} y \textbf{Productos}), la teor\'ia relacional exige la creaci\'on de una tabla asociativa con clave primaria compuesta:

\begin{lstlisting}[style=sqlstyle, caption={Implementaci\'on f\'isica de la tabla puente Order Details}]
-- Tabla asociativa que resuelve la relaci\'on N:M entre Orders y Products
CREATE TABLE order_details (
    order_id SMALLINT NOT NULL REFERENCES orders(order_id) ON DELETE CASCADE,
    product_id SMALLINT NOT NULL REFERENCES products(product_id),
    unit_price REAL NOT NULL CHECK (unit_price > 0),
    quantity SMALLINT NOT NULL CHECK (quantity > 0),
    discount REAL DEFAULT 0.0 CHECK (discount >= 0 AND discount <= 1),
    PRIMARY KEY (order_id, product_id)
);
\end{lstlisting}

---

\section{Taller de Consolidaci\'on del Cap\'itulo}

\begin{businessproblem}
Dise\~nar el esquema de base de datos relacional para un sistema de reservas de restaurantes gourmet, contemplando mesas, meseros asignados por turno, cartas de men\'u con platillos categorizados y comandas de mesas con m\'ultiples platos solicitados.
\end{businessproblem}

\subsection{Soluci\'on T\'ecnica Relacional}

\begin{lstlisting}[style=sqlstyle, caption={Esquema DDL completo para sistema de restaurantes}]
-- 1. Tablas Maestras
CREATE TABLE restaurant_tables (
    table_id SERIAL PRIMARY KEY,
    table_number VARCHAR(10) NOT NULL UNIQUE,
    seating_capacity SMALLINT NOT NULL CHECK (seating_capacity > 0),
    location_area VARCHAR(30) DEFAULT 'Salon Principal'
);

CREATE TABLE waitstaff (
    staff_id SERIAL PRIMARY KEY,
    first_name VARCHAR(50) NOT NULL,
    last_name VARCHAR(50) NOT NULL,
    shift_type VARCHAR(20) CHECK (shift_type IN ('Manana', 'Tarde', 'Noche'))
);

CREATE TABLE menu_dishes (
    dish_id SERIAL PRIMARY KEY,
    dish_name VARCHAR(100) NOT NULL,
    category VARCHAR(40) NOT NULL,
    price NUMERIC(8, 2) NOT NULL CHECK (price > 0)
);

-- 2. Entidad Transaccional y Detalle
CREATE TABLE restaurant_orders (
    order_id BIGSERIAL PRIMARY KEY,
    table_id INTEGER NOT NULL REFERENCES restaurant_tables(table_id),
    staff_id INTEGER NOT NULL REFERENCES waitstaff(staff_id),
    opened_at TIMESTAMPTZ DEFAULT CURRENT_TIMESTAMP,
    closed_at TIMESTAMPTZ,
    status VARCHAR(20) DEFAULT 'ABIERTA' CHECK (status IN ('ABIERTA', 'PAGADA', 'CANCELADA'))
);

CREATE TABLE order_dish_items (
    item_id BIGSERIAL PRIMARY KEY,
    order_id BIGINT NOT NULL REFERENCES restaurant_orders(order_id) ON DELETE CASCADE,
    dish_id INTEGER NOT NULL REFERENCES menu_dishes(dish_id),
    quantity SMALLINT NOT NULL DEFAULT 1 CHECK (quantity > 0),
    special_notes TEXT
);
\end{lstlisting}

\begin{engtip}
\begin{itemize}
    \item \textbf{Claves Primarias Compuestas vs Sustitutas:} En tablas intermedias con atributos propios (\textit{payload}) o con posibilidad de repetici\'on del mismo \'item en una comanda, es aconsejable emplear una clave sustituta (\texttt{BIGSERIAL}) manteniendo un \'indice \'unico en caso de requerir no-duplicidad.
\end{itemize}
\end{engtip}
